
% =========================================================
% 2. Motivation & Research Questions
% =========================================================
\section*{二、研究動機與研究問題}
\addcontentsline{toc}{section}{二、研究動機與研究問題}

\subsection*{(一) 研究動機}
本計畫聚焦於學術專題指導中「高度非線性議題演進」與「長週期進度追蹤」之複雜情境。在高等教育互動中，指導品質的核心取決於「教授人設穩定性（Persona Stability）」及其背後的學術權威立場。專題研究通常橫跨數月，涉及多次會議中研究假設之更迭、方法論之修正及實驗路徑之調整，對虛擬指導系統的記憶能力與角色穩定性提出嚴峻挑戰。

在記憶層面，現有大型語言模型（LLM）的記憶機制多侷限於片段式資訊存取，缺乏能表徵「研究過程決策因果依賴鏈」的結構化載體，難以在動態變化的研究決策中提供長程邏輯一貫的專業指導。MemoryBank 與 MemGPT 雖提供了長效存儲機制 \citep{a2305_10250v3_memorybank, a2310_08560v2_memgpt}，但面對高度邏輯嵌套的學術情境，仍難以達到圖式記憶（Graph-based Memory）在結構一致性與邏輯關係追蹤上的效果 \citep{a2602_05665v1_graphagentmemory}。

在角色一致性層面，有效的學術輔導要求虛擬導師維持批判性、方法論的嚴謹、對證據敏感並避免盲目迎合；系統若能忠實再現真實指導者的行為特徵，將有助於建立長期互動中的專業信任感。然而，LLM 的行為表現往往隨對話歷史而偏差，損害指導的可信度。近年 Persona Vectors 研究顯示，模型的人格特質可在激活空間（Activation Space）中以近似線性的方向被監測 \citep{a2507_21509v3_personavectors}；表徵工程（Representation Engineering）與推理時干預（Inference-time Steering）的進展則進一步證實，透過激活值加法（Activation Addition）等手段，可在不重新訓練的前提下，有效干預並校正模型的行為方向 \citep{turner-etal-2023-activation-addition, wehner-etal-2025-representation-engineering}。

基於上述背景，本研究的核心動機在於結合「記憶結構的穩固化」與「生成行為的即時校正」：透過「圖式長期記憶」維持跨會議的研究邏輯鏈，確保建議內容的事實一致性；同時利用「人設向量」於推理時進行即時監控與控制，確保指導風格的一致性。本計畫旨在建立一個在長週期專題輔導中更具邏輯記憶、可量測且指導風格穩定的虛擬教授系統。

\subsection*{(二) 研究問題}
本研究將以「長程一致性」為主軸，提出以下可驗證之研究問題（Research Questions, RQs）：

\begin{itemize}

  \item \textbf{RQ1（研究脈絡一致性）}：相較於僅使用文字記憶（非結構化摘要或片段檢索），以 Dependency Graph 建模之長期記憶是否能顯著降低跨會議對話中的研究邏輯斷裂與矛盾（例如目標變動、方法自相矛盾、TODO 遺失），並提升指導建議的可追溯性（traceability）與可檢核性？
  \item \textbf{RQ2（人設穩定性與漂移抑制）}：透過人設向量（Persona Vector）進行推理時干預，是否能在長程對話中維持指導教授互動策略的一致性，並以量化方式降低 persona drift 的發生程度？
  \item \textbf{RQ3（教育指導品質的外部效度）}：在真人盲測情境下，具備圖式長期記憶與人設向量監控之系統，是否能在指導品質量表（如 MCA-21）與 AI 導師評測框架中獲得顯著較佳的評分，並展現可部署的可靠性提升？
\end{itemize}
